\documentclass[11pt]{article}
\usepackage[margin=1in,top=.8in,bottom=.8in]{geometry}
\usepackage{amsmath,amssymb,amsthm,hyperref}
\hypersetup{colorlinks=true,linkcolor=blue,urlcolor=blue,citecolor=blue}
\setlength{\emergencystretch}{2em}
\newtheorem{theorem}{Theorem}
\newtheorem{lemma}[theorem]{Lemma}
\newcommand{\D}{\mathbb D}
\title{Unit-circle nonvanishing for rational Hartogs-triangle polynomials}
\author{Research proof and formalization packet}
\date{October 26, 2026}
\begin{document}
\maketitle
\begin{abstract}
We prove the full unit-circle nonvanishing assertion of Conjecture~4.20
of Edholm and Mathew. For every coprime pair of integers $m>n\geq1$,
their degree $2(m-n)$ palindromic polynomial $Q_{m,n}$ has no root on
the unit circle, and exactly $m-n$ roots in the punctured unit disc,
counting multiplicities. A finite coefficient reindexing expresses its
normalized boundary value as a sum of squared absolute values of floor
polynomials. Two adjacent floor polynomials have a difference consisting
of one monomial times $1-z$, which makes the sum strictly positive.
The source coefficient normalization and the full nonvanishing theorem
are verified in Lean. Earlier kernel formulas and the weighted
roots-of-unity identity of Almughrabi are explicitly credited.
\end{abstract}

\section{The source polynomial and the result}
Fix positive integers $m>n$ with $\gcd(m,n)=1$, and write $k=m-n$.
For an integer $u$ define the triangular coefficient
\begin{equation}\label{eq:D}
 D_m(u)=
 \begin{cases}
 u+1,&0\leq u\leq m-1,\\
 2m-1-u,&m\leq u\leq2m-2,\\
 0,&\text{otherwise}.
 \end{cases}
\end{equation}
These are the coefficients of $(1+w+\cdots+w^{m-1})^2$.
The numerator in Proposition~1.10, equations~(1.11)--(1.12), of the
published source \cite{EM} is the ordinary two-variable polynomial
\begin{equation}\label{eq:P}
 P_{m,n}(s,t)=\sum_{a=0}^{2m-2}\sum_{b=0}^{2n}
 D_m(a)D_m(mb+na+m+n-1-2mn)s^a t^b.
\end{equation}
The source defines
\begin{equation}\label{eq:Q}
 Q_{m,n}(z)=z^{1-2n}P_{m,n}(z,z).
\end{equation}
It proves that this is a polynomial of degree $2k$ with palindromic
coefficients. The apparent negative power in \eqref{eq:Q} cancels.
The formula \eqref{eq:P} is Almughrabi's earlier Proposition~2
\cite{Alm}; Edholm--Mathew also give an equivalent five-piece formula
with explicit ceiling functions in Theorem~1.3.

\begin{theorem}\label{thm:main}
For every coprime pair $m>n\geq1$ and every complex number $z$ with
$|z|=1$, one has $Q_{m,n}(z)\ne0$. Exactly $m-n$ roots of $Q_{m,n}$
lie in $\D^*=\{z\in\mathbb C:0<|z|<1\}$, counting multiplicities.
\end{theorem}

This is the full assertion of the published Conjecture~4.20
\cite{EM}, recorded as MathDB~368572. The source already establishes
differences $k=1,2,3,4,6$, and describes infinitely many explicit pairs
for every fixed difference. Those known cases are included in the
uniform theorem and receive no separate discovery claim.

\section{A finite sum of squares}
For each integer $0\leq r<m$, define
\begin{equation}\label{eq:A}
 A_r(z)=\sum_{p=0}^{m-1}z^{e_r(p)},\qquad
 e_r(p)=p-\left\lfloor\frac{np+r}{m}\right\rfloor.
\end{equation}
The exponents satisfy $0\leq e_r(p)\leq k$. Indeed,
$np+r<m(p+1)$ gives the lower bound, and
$e_r(p)=\lceil(kp-r)/m\rceil\leq k$ gives the upper bound.

\begin{lemma}\label{lem:SOS}
For every nonzero complex $z$,
\begin{equation}\label{eq:LaurentSOS}
 Q_{m,n}(z)=z^k\sum_{r=0}^{m-1}A_r(z)A_r(z^{-1}).
\end{equation}
Consequently, if $|z|=1$, then
\begin{equation}\label{eq:SOS}
 z^{-k}Q_{m,n}(z)=\sum_{r=0}^{m-1}|A_r(z)|^2.
\end{equation}
\end{lemma}

\begin{proof}
We give a direct coefficient proof, keeping the source normalization.
For $a,c\in\{0,\ldots,2m-2\}$, the substitution
\[
 c=mb+na+m+n-1-2mn
\]
identifies the nonzero terms of \eqref{eq:P} with the congruence
$c-na\equiv n-1\pmod m$. Conversely that congruence determines the
integer $b$, and $0\leq b\leq2n$: the numerator determining $mb$
is strictly between $-m$ and $m(2n+1)$, and is divisible by $m$.
Thus \eqref{eq:Q} becomes
\begin{equation}\label{eq:compressed}
 Q_{m,n}(z)=
 \sum_{\substack{0\leq a,c\leq2m-2\\c-na\equiv n-1\ (\mathrm{mod}\ m)}}
 D_m(a)D_m(c)\,
 z^{(ka+c-n+1)/m}.
\end{equation}
The exponent is an integer; the same substitution shows that it equals
$a+b+1-2n$.

The integer $D_m(a)$ counts the ordered pairs $p,p_2\in\{0,\ldots,m-1\}$
with $p+p_2=a$. Complementing the second member writes
$a=p+m-1-p'$ with $0\leq p,p'<m$. Similarly write
$c=q+m-1-q'$ with $0\leq q,q'<m$. The congruence in
\eqref{eq:compressed} becomes
\begin{equation}\label{eq:residue}
 q-np\equiv q'-np'\pmod m.
\end{equation}
Group these pairs by their common residue $r\in\{0,\ldots,m-1\}$.
For a fixed $p$ and $r$, there is exactly one $q\in\{0,\ldots,m-1\}$,
namely
\[
 q=np+r-m\left\lfloor\frac{np+r}{m}\right\rfloor.
\]
The analogous formula holds for $q'$. The exponent in
\eqref{eq:compressed} now equals
\[
 \frac{k(p+m-1-p')+q+m-1-q'-n+1}{m}
 =k+e_r(p)-e_r(p').
\]
All changes of variables above are bijections of finite index sets,
including their multiplicities $D_m(a)D_m(c)$. Summing first over
$p,p'$ and then over $r$ proves \eqref{eq:LaurentSOS}.
For unit $z$, the real coefficients of $A_r$ give
$A_r(z^{-1})=\overline{A_r(z)}$, proving \eqref{eq:SOS}.
\end{proof}

\section{Strict positivity and root allocation}
\begin{lemma}\label{lem:strict}
The two polynomials $A_0$ and $A_1$ have no common nonzero complex root.
\end{lemma}
\begin{proof}
Coprimality gives a unique $p_*\in\{1,\ldots,m-1\}$ with
$np_*\equiv m-1\pmod m$. Write
$np_*=mh+m-1$. In passing from $r=0$ to $r=1$ in
\eqref{eq:A}, the floor increases only for this one value of $p$.
Consequently
\begin{equation}\label{eq:difference}
 A_1(z)-A_0(z)=z^{p_*-h-1}(1-z).
\end{equation}
The exponent is nonnegative: the general exponent bound for $A_1$
gives $h+1\leq p_*$. If $z\ne0,1$, the right side of
\eqref{eq:difference} is nonzero. At $z=1$, one has $A_0(1)=m\ne0$.
\end{proof}

\begin{proof}[Proof of Theorem~\ref{thm:main}]
On the unit circle every summand in \eqref{eq:SOS} is nonnegative,
and Lemma~\ref{lem:strict} makes at least one of the first two summands
strictly positive. Hence $z^{-k}Q_{m,n}(z)>0$ and $Q_{m,n}(z)\ne0$.

For completeness, the same coefficient identity recovers the algebra
needed to count roots. Because $0\leq e_r(p)\leq k$, every term on the
right of \eqref{eq:LaurentSOS} has exponent in $[0,2k]$. The expression
is an ordinary polynomial with nonnegative integer coefficients, and
swapping $p$ and $p'$ shows that it is palindromic. In $A_0$, the
indices $p=0$ and $p=m-1$ give exponents $0$ and $k$, respectively.
Thus both the constant and degree-$2k$ coefficients of $Q_{m,n}$ are
positive. The degree is exactly $2k$ and zero is not a root.
Reciprocity pairs each root $\zeta$ with $\zeta^{-1}$ with the same
multiplicity. Since there is no unit-circle root, inversion gives a
multiplicity-preserving bijection between roots inside and outside the
unit disc. The fundamental theorem of algebra counts $2k$ roots in
all, so exactly $k$ lie in $\D^*$.
\end{proof}

\section{Relation to the older kernel formulas}
The finite roots-of-unity machinery underlying this calculation is
classical. In particular, the weighted geometric-square identity itself
already appears in Almughrabi \cite[eq.~(24)]{Alm}; it is not a new
claim of this paper. To make the relation precise, put
$S_m(u)=1+u+\cdots+u^{m-1}$ and let $\omega$ be a primitive $m$th
root of unity. For $b\ne0$, root-of-unity filtering gives
\begin{align}\label{eq:oldfilter}
 &b^{m+n-1-2mn}P_{m,n}(ab^n,b^m)\notag\\
 &\hspace{5mm}=\frac1m\sum_{j=0}^{m-1}\omega^{(1-n)j}
 S_m(a\omega^{-nj})^2 S_m(b\omega^j)^2.
\end{align}
The identity is Almughrabi's formula after changing $j$ to $-j$ and
matching its normalization. For unit $u$,
$S_m(u)^2=u^{m-1}|S_m(u)|^2$. The phases in each summand of
\eqref{eq:oldfilter} therefore cancel, yielding a nonnegative sum.
The direct finite reindexing in Lemma~\ref{lem:SOS} gives the explicit
floor polynomials and the strictness certificate
\eqref{eq:difference}, including all resonant roots of unity.

Chakrabarti--Konkel--Mainkar--Miller \cite{CKMM} supply the earlier
triangular-coefficient kernel formula, and Almughrabi \cite{Alm}
simplifies it in two complex dimensions. These formulas receive full
prior credit. Neither the actual earlier formula papers nor the
June~2026 source examined in this review asserts the full boundary
positivity conclusion proved here. The source explicitly retains
Conjecture~4.20 and mentions a further paper in preparation; no public
full version of that successor was located. This bounded comparison
supports the proposed source-conjecture resolution and does not certify
absolute worldwide priority.

\section{The complex-analytic consequence and formal scope}
Edholm--Mathew's source formula is
\[
 K_{m/n}(Z,W)=\frac{P_{m,n}(s,t)}{m\pi^2(1-t)^2(t^n-s^m)^2},
 \qquad s=Z_1\overline{W_1},\quad t=Z_2\overline{W_2},
\]
on $\mathbb H_{m/n}=\{(Z_1,Z_2):|Z_1|^{m/n}<|Z_2|<1\}$.
The theorem supplies roots $z\in\D^*$ of $P_{m,n}(z,z)$.
The diagonal $(z,z)$ belongs to $\mathbb H_{m/n}$, and the denominator
there is nonzero. Such a pair of products is realized by taking
$Z=(\sqrt{|z|}e^{i\arg z},\sqrt{|z|}e^{i\arg z})$ and
$W=(\sqrt{|z|},\sqrt{|z|})$, both in the domain.
Thus the source kernel formula gives zeros for every rational parameter
$m/n>1$. This is the source's stated consequence of its polynomial
conjecture, not a separate counted result. The passage uses the
published kernel formula; the Lean certificate concerns the explicit
polynomials and finite algebra, not the theory constructing Bergman
kernels or proper holomorphic maps.

The ten-module Lean certificate verifies the signed source coefficient
$D_m$, its support and literal double-sum normalization, the affine
coefficient reindexing, triangular convolution, complementary-pair and
residue reindexings, the exact unit-circle sum-of-squares identity, the
modular carry and strictness argument, and nonvanishing for every
coprime pair. It also constructs an ordinary polynomial over the
natural numbers, maps it into $\mathbb C[X]$, proves equality with the
literal source expression at every nonzero argument, bounds every
selected integer exponent before conversion to a natural exponent,
and proves exact degree $2k$, nonzero constant coefficient and
reciprocity. Full complex factorization proves the reciprocal-root
multiset identity. The final source endpoint therefore certifies both
unit-circle nonvanishing and exactly $k$ punctured-disc roots,
including all multiplicities. No source, SOS, degree, reciprocity or
root-count premise is left unproved.

The main endpoint is \path{source_full_conjecture_368572} in namespace
\texttt{BergmanUnitCircle}. Fresh pinned replay checks all 78 local
declarations with warnings treated as errors. Axiom inspection permits
only \texttt{propext}, \texttt{Classical.choice}, and
\texttt{Quot.sound}. The compiler is Lean~4.34.1; exact compiler and
mathlib commits, respectively, are
\begin{center}\small\ttfamily
5045d0056413266e57c625dcd7c365b10e377c52\\
d13f23b723b8a846827a245b89c10fc7d3f11612.
\end{center}
The packet contains exact dependency pins, independent reviews,
all-declaration replay receipts and types/axioms. No numerical
experiment, proof hole, added axiom, unsafe declaration or native
computation assumption is used. The general Laurent identity above
is proved mathematically; the formally certified SOS specializes it
to unit inputs, which suffices for the full stated source conclusion.

For a separate arithmetic cross-check, the literal ceiling-function
five-piece formula in the source was constructed with exact integer
polynomials for every coprime pair $2\leq m\leq80$. All $1,965$ cases
agree with \eqref{eq:LaurentSOS} and \eqref{eq:difference}. These checks
corroborate the definitions and are not assumptions in the uniform
proof.

\begin{thebibliography}{9}
\bibitem{EM}
L.~D.~Edholm and V.~T.~Mathew,
\emph{Arithmetic properties and zeros of the Bergman kernel on a class
of quotient domains},
Mathematische Zeitschrift \textbf{313} (2026), article~43,
\href{https://doi.org/10.1007/s00209-026-04059-8}{doi:10.1007/s00209-026-04059-8};
also \href{https://arxiv.org/abs/2505.20489v3}{arXiv:2505.20489v3}.

\bibitem{Alm}
R.~Almughrabi,
\emph{Bergman kernels of two dimensional monomial polyhedra},
Complex Analysis and Operator Theory \textbf{17} (2023), article~103,
\href{https://doi.org/10.1007/s11785-023-01401-4}{doi:10.1007/s11785-023-01401-4};
also \href{https://arxiv.org/abs/2303.14268}{arXiv:2303.14268}.

\bibitem{CKMM}
D.~Chakrabarti, A.~Konkel, M.~Mainkar and E.~Miller,
\emph{Bergman kernels of elementary Reinhardt domains},
Pacific Journal of Mathematics \textbf{306} (2020), 67--93;
\href{https://arxiv.org/abs/1909.03164}{arXiv:1909.03164}.
\end{thebibliography}
\end{document}
